\section{Implementation Details}
\label{appendix:hyperparams}

Below, we explain the implementation details for CURL in the DMControl setting. Specifically, we use the SAC algorithm as the RL objective coupled with CURL and build on top of the publicly released implementation from \citet{yarats2019improving}. We present in detail the hyperparameters for the architecture and optimization. We do not use any extra hyperparameter for balancing the contrastive loss and the reinforcement learning losses. Both the objectives are weighed equally in the gradient updates.

\begin{table}[h]
\caption{Hyperparameters used for DMControl CURL experiments. Most hyperparameters values are unchanged across environments with the exception for action repeat, learning rate, and batch size.}

\label{table:hyperparameters}
\vskip 0.15in
\begin{center}
\begin{small}
\begin{tabular}{ll}
\toprule
\textbf{Hyperparameter} & \textbf{Value}  \\
\midrule
Random crop    & True  \\ 
Observation rendering    & $(100,100)$  \\ 
Observation downsampling    & $(84,84)$  \\ 
Replay buffer size    & $100000$ \\ 
Initial steps    & $1000$  \\ 
Stacked frames    & $3$  \\ 
Action repeat    & $2$ finger, spin; walker, walk\\
 & $8$ cartpole, swingup \\
 & $4$ otherwise  \\
Hidden units (MLP)    & $1024$  \\ 
Evaluation episodes    & $10$  \\ 
Optimizer    & Adam  \\ 
$(\beta_1,\beta_2) \rightarrow (f_\theta, \pi_\psi, Q_\phi)$   & $(.9,.999)$  \\
$(\beta_1,\beta_2) \rightarrow (\alpha)$   & $(.5,.999)$  \\
Learning rate $(f_\theta, \pi_\psi, Q_\phi)$     & $2e-4$ cheetah, run \\
& $1e-3$ otherwise\\ 
Learning rate ($\alpha$) & $1e-4$ \\

Batch Size    & $512$  \\ 
$Q$ function EMA $\tau$ & $0.01$ \\
Critic target update freq & $2$ \\
Convolutional layers & $4$ \\
Number of filters & $32$ \\
Non-linearity & ReLU \\
Encoder EMA $\tau$ & $0.05$ \\
Latent dimension & $50$ \\
Discount $\gamma$ & $.99$ \\
Initial temperature & $0.1$ \\


\bottomrule
\end{tabular}
\end{small}
\end{center}
\vskip -0.1in
\end{table}

{\textbf{Architecture:}} We use an encoder architecture that is similar to \cite{yarats2019improving}, which we sketch in PyTorch-like pseuodocode below. The actor and critic both use the same encoder to embed image observations. A full list of hyperparameters is displayed in Table \ref{table:hyperparameters}.

For contrastive learning, CURL utilizes momentum for the key encoder \cite{kaiming2019moco} and a bi-linear inner product as the similarity measure \cite{oord2018representation}. Performance curves ablating these two architectural choices are shown in Figure \ref{fig:method_ablations}.

\begin{figure}[H]
\vskip 0.2in
\begin{center}
\centerline{\includegraphics[width=\columnwidth]{plots/reco_inner_prod_and_ema.pdf}}
\caption{Performance on cheetah-run environment ablated two-ways: (left) using the query encoder or exponentially moving average of the query encoder for encoding keys (right) using the bi-linear inner product as in \cite{oord2018representation} or the cosine inner product as in \citet{kaiming2019moco,chen2020simclr}}
\label{fig:method_ablations}
\end{center}
\vskip -0.2in
\end{figure}

Pseudo-code for the architecture is provided below:

\lstinputlisting[language=Python]{scripts/encode.py}

{\textbf{Terminology}:} A common point of confusion is the meaning  ``training steps." We use the term \textit{environment steps} to denote the amount of times the simulator environment is stepped through and \textit{interaction steps} to denote the number of times the agent steps through its policy. The terms \textit{action repeat} or \textit{frame skip} refer to the number of times an action is repeated when it's drawn from the agent's policy. For example, if action repeat is set to 4, then 100k interaction steps is equivalent to 400k environment steps.


{\textbf{Batch Updates:}} After initializing the replay buffer with observations extracted by a random agent, we sample a batch of observations, compute the CURL objectives, and step through the optimizer. Note that since queries and keys are generated by data-augmenting an observation, we can generate arbitrarily many keys to increase the contrastive batch size without sampling any additional observations.

{\textbf{Shared Representations:}}  
The objective of performing contrastive learning together with RL is to ensure that the shared encoder learns rich features that facilitate sample efficient control. There is a subtle coincidental connection between MoCo and off-policy RL. Both the frameworks adopt the usage of a momentum averaged (EMA) version of the underlying model. In MoCo, the EMA encoder is used for encoding the keys (targets) while in off-policy RL, the EMA version of the Q-networks are used as targets in the Bellman error \cite{mnih2015human, haarnoja2018soft}. Thanks to this connection, CURL shares the convolutional encoder, momentum coefficient and EMA update between contrastive and reinforcement learning updates for the shared parameters. The MLP part of the critic that operates on top of these convolutional features has a separate momentum coefficient and update decoupled from the image encoder parameters.

{\textbf{Balancing Contrastive and RL Updates}:}
While past work has learned hyperparameters to balance the auxiliary loss coefficient or learning rate relative to the RL objective \cite{jaderberg2016reinforcement, yarats2019improving}, CURL does not need any such adjustments. We use both the contrastive and RL objectives together with equal weight and learning rate. This simplifies the training process compared to other methods, such as training a VAE jointly \cite{hafner2018learning,hafner2019dream,lee2019stochastic}, that require careful tuning of coefficients for representation learning.


{\textbf{Differences in Data Collection between Computer Vision and RL Settings}:} 
There are two key differences between contrastive learning in the computer vision and RL settings because of their different goals. Unsupervised feature learning methods built for downstream vision tasks like image classification assume a setting where there is a large static dataset of unlabeled images.  On the other hand, in RL, the dataset changes over time to account for the agent's new experiences. Secondly, the size of the memory bank of labeled images and dataset of unlabeled ones in vision-based settings are 65K and 1M (or 1B) respectively. The goal in vision-based methods is to learn from millions of unlabeled images. On the other hand, the goal in CURL is to develop sample-efficient RL algorithms. For example, to be able to solve a task within 100K timesteps (approximately 2 hours in real-time), an agent can only ingest 100K image frames.
 
 Therefore, unlike MoCo, CURL does not use a memory bank for contrastive learning. Instead, the negatives are constructed on the fly for every minibatch sampled from the agent's replay buffer for an RL update similar to SimCLR. The exact implementation is provided as a PyTorch-like code snippet in \ref{pytorch}.


{\textbf{Data Augmentation:}}  


Random crop data augmentation has been crucial for the performance of deep learning based computer vision systems in object recognition, detection and segmentation \cite{krizhevsky2012, szegedy2015, cubuk2019, chen2020simclr}. However, similar augmentation methods have not seen much adoption in the field of RL even though several benchmarks use raw pixels as inputs to the model.

CURL adopts the random crop data augmentation as the stochastic data augmentation applied to a frame stack. To make it easier for the model to correlate spatio-temporal patterns in the input, we apply the same random crop (in terms of box coordinates) across all four frames in the stack as opposed to extracting different random crop positions from each frame in the stack. Further, unlike in computer vision systems where the aspect ratio for random crop is allowed to be as low as 0.08, we preserve much of the spatial information as possible and use a constant aspect ratio of 0.84 between the original and cropped. In our experiments, data augmented samples for CURL are formed by cropping $84\times84$ frames from an input frame of $100\times100$. 

{\textbf{DMControl:}} We render observations at $100 \times 100$ and randomly crop $84 \times 84$ frames. For evaluation, we render observations at $100 \times 100$ and center crop to $84 \times 84$ pixels. We found that implementing random crop efficiently was extremely important to the success of the algorithm. We provide pseudocode below:

\lstinputlisting[language=Python]{scripts/crop.py}


\section{Atari100k Implementation Details}

The flexibility of CURL allows us to apply it to discrete control setting with minimal modifications. Similar to our rationale for picking SAC as the baseline RL algorithm to couple CURL with (for continuous control), we pick the data-efficient version of Rainbow DQN (Efficient Rainbow) \cite{van2019use} for Atari100K which performs competitively with an older version of SimPLe (most recent version has improved numbers). In order to understand {\it specifically} what the gains from CURL are without any other changes, we adopt the {\it exact} same hyperparameters specified in the paper~\cite{van2019use} (including a modified convolutional encoder that uses larger kernel size and stride of 5). We present the details in Table \ref{table:hyperparametersatari}. Similar to DMControl, the contrastive objective and the RL objective are weighted equally for learning (except for Pong, Freeway, Boxing and PrivateEye for which we used a coefficient of $0.05$ for the momentum contastive loss. On a large majority (22 out of 26) of the games, we do not use this adjustment. While it is standard practice to use the same hyperparameters for all games in Atari, papers proposing auxiliary losses have adopted a different practice of using game specific coefficients~\cite{jaderberg2016reinforcement}.). We use the Efficient Rainbow codebase from \url{https://github.com/Kaixhin/Rainbow} which has a reproduced version of \citet{van2019use}. We evaluate with $20$ random seeds and report the mean score for each game given the high variance nature of the Atari100k steps benchmark. We restrict ourselves to using grayscale renderings of image observations and use random crop of frame stack as data augmentation.

\begin{table}[h]
\caption{Hyperparameters used for Atari100K CURL experiments. Hyperparameters are unchanged across games.}

\label{table:hyperparametersatari}
\vskip 0.15in
\begin{center}
\begin{small}
\begin{tabular}{ll}
\toprule
\textbf{Hyperparameter} & \textbf{Value}  \\
\midrule
Random crop    & True  \\ 
Image size    & $(84,84)$  \\ 
Data Augmentation & Random Crop (Train) \\
Replay buffer size    & $100000$ \\ 
Training frames & $400000$ \\ 
Training steps & $100000$ \\
Frame skip & $4$ \\
Stacked frames    & $4$  \\ 
Action repeat    & $4$ \\
Replay period every & $1$ \\
Q network: channels & $32$, $64$ \\
Q network: filter size & $5\times 5, 5\times 5$ \\
Q network: stride & $5$, $5$ \\
Q network: hidden units & $256$ \\
Momentum (EMA for CURL) $\tau$ & $0.001$  \\
Non-linearity & ReLU \\
Reward Clipping   & $[-1, 1]$  \\ 
Multi step return & $20$ \\
Minimum replay size for sampling & $1600$ \\ 
Max frames per episode & $108$K \\
Update & Distributional Double Q \\
Target Network Update Period & every $2000$ updates \\
Support-of-Q-distribution & $51$ bins \\ 
Discount $\gamma$ & $0.99$ \\
Batch Size & $32$  \\
Optimizer    & Adam  \\ 
Optimizer: learning rate & $0.0001$ \\
Optimizer: $\beta1$ & $0.9$ \\
Optimizer: $\beta2$ & $0.999$ \\ 
Optimizer $\epsilon$ & $0.000015$ \\
Max gradient norm & $10$ \\ 
Exploration & Noisy Nets   \\
Noisy nets parameter & $0.1$ \\
Priority exponent & $0.5$ \\
Priority correction & $0.4 \rightarrow 1$ \\
Hardware & CPU \\
\bottomrule
\end{tabular}
\end{small}
\end{center}
\vskip -0.1in
\end{table}


\section{Benchmarking Data Efficiency}

Tables \ref{table:500kscores} and \ref{table:100katariscores} show the episode returns of DMControl100k, DMControl500k, and Atari100k across CURL and a number of pixel-based baselines. CURL outperforms all baseline pixel-based methods across experiments on both DMControl100k and DMControl500k. On Atari100k experiments, CURL coupled with Eff Rainbow outperforms the baseline on the majority of games tested (19 out of 26 games).

 \begin{figure}[h]
 \begin{center}
   \centerline{\includegraphics[width=\columnwidth]{plots/reco_dreamer_data_efficiency.pdf}}
   \caption{The number of steps it takes a prior leading pixel-based method, Dreamer, to achieve the same score that CURL achieves at 100k training steps (clipped at 1M steps). On average, CURL is 4.5x more data-efficient. We chose Dreamer because the authors \cite{hafner2019dream} report performance for all of the above environments while other baselines like SLAC and SAC+AE only benchmark on 4 and 6 environments, respectively. For further comparison of CURL with these methods, the reader is referred to Table \ref{table:500kscores} and Figure \ref{fig:curl_main_result}.
   }
   \label{fig:curl_sample_efficiency}
 \end{center}
 \vskip -0.2in
 \end{figure}
 


\section{Further Investigation of Data-Efficiency in Contrastive RL}

To further benchmark CURL's sample-efficiency, we compare it to state-based SAC on a total of 16 DMControl environments. Shown in Figure \ref{fig:curl_all_dmc}, CURL matches state-based data-efficiency on most of the environments, but lags behind state-based SAC on more challenging environments.



\begin{figure}[!ht]
 \begin{center}
   \centerline{\includegraphics[width=11cm]{plots/reco_dmc_state_16.pdf}}
   \caption{CURL compared to state-based SAC run for 3 seeds on each of 16 selected DMControl environments. For the 6 environments in \ref{fig:curl_main_result}, CURL performance is averaged over 10 seeds.}
   \label{fig:curl_all_dmc}
 \end{center}
 \end{figure}
 

\section{Ablations}\label{appendix:ablation}

\subsection{Learning Temporal Dynamics}
\begin{figure}[t]
 \begin{center}
   \centerline{\includegraphics[width=\columnwidth]{plots/reco_temporal_ablation.pdf}}
   \caption{CURL with temporal and visual discrimination (red) compared to CURL with only visual discrimination (green). In most settings, the variant with temporal variant outperforms the purely visual variant of CURL. The two exceptions are reacher and ball in cup environments, suggesting that learning dynamics is not necessary for those two environments. Note that the walker environment was run with action repeat of 4, whereas walker walk in the main results Table \ref{table:500kscores} and Figure \ref{fig:curl_all_dmc} was run with action repeat of 2. }
   \vspace{-6mm}
   \label{fig:curl_temporal_ablation}
 \end{center}
 \end{figure}
 
 To gain insight as to whether CURL learns temporal dynamics across the stacked frames, we also train a variant of CURL where the discriminants are individual frames as opposed to stacked ones. This can be done by sampling stacked frames from the replay buffer but only using the first frame to update the contrastive loss:

\begin{lstlisting}[language=python]
f_q = x_q[:,:3,...] # (B,C,H,W), C=9.
f_k = x_k[:,:3,...]
\end{lstlisting}


During the actor-critic update, frames in the batch are encoded individually into latent codes, which are then concatenated before being passed to a dense network.

\begin{lstlisting}[language=python]
# x: (B,C,H,W), C=9.
z1 = encode(x[:,:3,...])
z2 = encode(x[:,3:6,...])
z3 = encode(x[:,6:9,...])
z = torch.cat([z1,z2,z3],-1)
\end{lstlisting}


Encoding each frame indiviudally ensures that the contrastive objective only has access to visual discriminants. Comparing the visual and spatiotemporal variants of CURL in Figure \ref{fig:curl_temporal_ablation} shows that the variant trained on stacked frames outperforms the visual-only version in most environments. The only exceptions are reacher and ball-in-cup environments. Indeed, in those environments the visual signal is strong enough to solve the task optimally, whereas in other environments, such as walker and cheetah, where balance or coordination is required, visual information alone is insufficient.

\subsection{Increasing Gradient Updates per Agent Step}

Although most baselines we benchmark against use one gradient update per agent step, it was recently empirically shown that increasing the ratio of gradients per step improves data-efficiency in RL \cite{kielak2020rainbow}. This finding is also supported by SLAC \cite{lee2019stochastic}, where results are shown with a ratio of 1:1 (SLACv1) and 3:1 (SLACv2). We 


\begin{table}[h!]
\caption{Scores achieved by CURL and SLAC when run with a 3:1 ratio of gradient updates per agent step on DMControl500k and DMControl100k. CURL achieves state-of-the-art performance on the majority (\textbf{3} out of \textbf{4}) environments on DMControl500k. Performance of both algorithms is improved relative to the 1:1 ratio reported for all baselines in Table \ref{table:500kscores} but at the cost of significant compute and wall-clock time overhead. }
\label{table:three_grad}
\vskip 0.15in
\begin{center}
\begin{small}
\begin{sc}
\begin{tabular}{lcc}
\toprule
DMControl500k & CURL & SLACv2 \\
\midrule
Finger, spin    & \textbf{923 $\pm$ 50} & 884 $\pm$ 98 \\
Walker, walk    & \textbf{911 $\pm$ 35} & 891  $\pm$ 60\\\
Cheetah, run    & 545 $\pm$ 39 & \textbf{791 $\pm$ 37}  \\
Ball in cup, catch   & \textbf{948 
$\pm$ 21 }  & 885 $\pm$ 154 \\
\midrule
DMControl100k & CURL & SLACv2 \\
\midrule
Finger, spin    & \textbf{ 741$\pm$ 118 }  & 728 $\pm$212  \\
Walker, walk    & 428 $\pm$ 59 & \textbf{513 $\pm$ 41}  \\\
Cheetah, run    & 314 $\pm$ 46 & \textbf{438 $\pm$ 76}  \\
Ball in cup, catch    & \textbf{ 899 $\pm$ 47 }& 837 $\pm$ 147  \\
\bottomrule
\end{tabular}
\end{sc}
\end{small}
\end{center}
\vskip -0.1in
\end{table}

 \subsection{Decoupling Representation Learning from Reinforcement Learning}

Typically, Deep RL representations depend almost entirely on the reward function specific to a task. However, hand-crafted representations such as the proprioceptive state are independent of the reward function. It is much more desirable to learn reward-agnostic representations, so that the same representation can be re-used across different RL tasks. We test whether CURL can learn such representations by comparing CURL to a variant where the critic gradients are backpropagated through the critic and contrastive dense feedforward networks but stopped before reaching the convolutional neural network (CNN) part of the encoder. 

Scores displayed in Figure \ref{fig:detached_plots} show that for many environments, the detached CNN representations are sufficient to learn an optimal policy. The major exception is the cheetah environment, where the detached representation significantly under-performs. Though promising, we leave further exploration of task-agnostic representations for future work.

  \begin{figure}[!ht]
 \begin{center}
   \centerline{\includegraphics[width=\columnwidth]{plots/reco_detached_ablation.pdf}}
   \caption{CURL where the CNN part of the encoder receives gradients from both the contrastive loss and critic (red) compared to CURL with the convolutional part of the encoder trained only with the contrastive objective (green). The detached encoder variant is able to learn representations that enable near-optimal learning on most environments, except for cheetah. As in Figure \ref{fig:curl_temporal_ablation}, the walker environment was run with action repeat of 4.}
   \label{fig:detached_plots}
 \end{center}
  \vskip -.2in
  \vspace{-6mm}
 \end{figure}

 \subsection{Removing Data Augmentation for the Actor Critic}
 \label{noaug}

Our main results involve the use of data augmentations to regularize both the contrastive and SAC objectives. Here, we investigate whether the contrastive representations alone are sufficient for learning effective policies. In these experiments, we only augment the data for the contrastive objective but not for the SAC agent. As a result, data augmentation is used only to learn features but does not influence the control policy. The pseudocode is shown below:

\lstinputlisting[language=Python]{scripts/no_aug.py}


\begin{figure}[H]
\vskip 0.2in
\begin{center}
\centerline{\includegraphics[width=.7\columnwidth]{plots/curl_no_aug_sac.png}}
\caption{CURL with no data augmentations passed to the SAC agent improves the performance of the baseline pixel SAC by a mean of 2.0x / median of 1.7x on DMControl500k. For these runs we use a smaller batch size of 128 than the 512 batch size used for results in Table \ref{fig:curl_main_result}. While the constastive loss alone improves over the pixel SAC baseline, most environments benefit from data augmentation also being passed to the SAC agent.}
\label{fig:curl_noaug}
\end{center}
\vskip -0.2in
\end{figure}

DMControl500k results plotted in Figure~\ref{fig:curl_noaug} show that, on average, features learned through the contrastive loss alone improve the pixel SAC baseline by 2x. Augmenting the input passed to the SAC algorithm further improves performance.

\subsection{Predicting State from Pixels}

Despite improved sample-efficiency on most DMControl tasks, there is still a visible gap between the performance of SAC on state and SAC with CURL in some environments. Since CURL learns representations by performing instance discrimination across stacks of three frames, it's possible that the reason for degraded sample-efficiency on more challenging tasks is due to partial-observability of the ground truth state. 

To test this hypothesis, we perform supervised regression $(X,Y)$ from pixels $X$ to the proprioceptive state $Y$, where each data point $x \in X$ is a stack of three consecutive frames and $y \in Y$ is the corresponding state extracted from the simulator. We find that the error in predicting the state from pixels correlates with the policy performance of pixel-based methods. Test-time error rates displayed in Figure \ref{fig:statereg} show that environments that CURL solves as efficiently as state-based SAC have low error-rates in predicting the state from stacks of pixels. The prediction error increases for more challenging environments, such as cheetah-run and walker-walk. Finally, the error is highest for environments where current pixel-based methods, CURL included, make no progress at all \cite{tassa2018deepmind}, such as humanoid and swimmer.

This investigation suggests that degraded policy performance on challenging tasks may result from the lack of requisite information about the underlying state in the pixel data used for learning representations. We leave further investigation for future work.

\begin{figure}[!ht]
\begin{center}
\centerline{\includegraphics[width=7cm]{plots/state_regs.png}}
\caption{Test-time mean squared error for predicting the proprioceptive state from pixels on a number of DMControl environments. In DMControl, environments fall into two groups - where the state corresponds to either (a) positions and velocities of the robot joints or (b) the joint angles and angular velocities.}
   \label{fig:statereg}
\end{center}
\vskip -0.2in
\end{figure}

\subsection{CURL + Efficient Rainbow Atari runs}

We report the scores (Tables \ref{tab:atari1} and \ref{tab:atari2}) for 20 seeds across the 26 Atari games in the Atari100k benchmark for CURL coupled with Efficient Rainbow. The variance across multiple seeds is considerably high in this benchmark. Therefore, we report the scores for each of the seeds along with the mean and standard deviation for each game. 

\begin{table*}[!ht]%
\centering
\small
\setlength\tabcolsep{2.5pt}
\begin{tabular}{c|c|c|c|c|c|c|c|c|c|c|c|c}
\hline
Pacman &	Frostbite &	Asterix  &	KungFuMaster &	Kangaroo &	Gopher &	RoadRunner &	JamesBond &	BattleZone &	Seaquest &	Assault &	Krull &	Qbert			 \\
\hline
1287 &	2292 &	850 &	8470 &	600 &	1036 &	2820 &	305 &	18100 &	322 &	634.2 &	3404.3 &	1020 \\
1608 &	1046 &	525 &	10870 &	2280 &	574 &	3190 &	265 &	18200 &	236 &	696.8 &	2443.5 &	650  \\
1466 &	1209 &	655 &	10920 &	1940 &	540 &	7840 &	335 &	26800 &	352 &	655.2 &	6791.4 &	830  \\
1430 &	255 &	565 &	7730 &	1140 &	618 &	12060 &	145 &	21300 &	386 &	443 &	3022.5 &	902.5  \\
1114 &	426 &	715 &	17525 &	520 &	534 &	8340 &	565 &	7900 &	458 &	546 &	3892.2 &	3957.5  \\
1083 &	2280 &	715 &	3560 &	600 &	596 &	6920 &	565 &	8100 &	224 &	564.9 &	3505.5 &	772.5 \\
2301 &	259 &	770 &	10940 &	600 &	502 &	2230 &	350 &	12000 &	282 &	514.4 &	2564.1 &	782.5  \\
1128 &	335 &	980 &	23420 &	900 &	998 &	4250 &	365 &	16500 &	339 &	516.6 &	4079.7 &	727.5  \\
1184 &	1409 &	665 &	15160 &	600 &	950 &	1570 &	140 &	23900 &	526 &	661.5 &	2376.4 &	705 \\
1510 &	258 &	610 &	15370 &	730 &	544 &	6300 &	425 &	19900 &	436 &	664.5 &	4161.8 &	757.5  \\
2343 &	335 &	905 &	22260 &	600 &	796 &	3100 &	315 &	10000 &	272 &	529 &	3311.1 &	647.5  \\
1063 &	1062 &	800 &	17320 &	880 &	522 &	1060 &	335 &	11200 &	428 &	445.2 &	2517.3 &	562.5  \\
2040 &	1542 &	675 &	31820 &	220 &	392 &	6050 &	735 &	9700 &	358 &	573.3 &	3764.7 &	2425  \\
1195 &	1102 &	795 &	23360 &	920 &	780 &	11810 &	950 &	23500 &	533 &	531.3 &	10150.2 &	1112.5  \\
1343 &	2461 &	585	 & 27460 & 600 &	792 &	4630 &	520 &	10500 &	968 &	663.6 &	2883.6 &	527.5  \\
1354 &	257 &	865 &	7770 &	2300 &	454 &	2530 &	755 &	18100 &	314 &	795.3 &	5123.7	 & 472.5  \\
1925 &	513 &	730 &	8820 &	320 &	564 &	6840 &	750 &	9000 &	378 &	633 &	3652.5 &	610  \\
1228 &	1826 &	680 &	2980 &	600 &	522 &	6580 &	795 &	8900 &	168 &	674.1 &	2376.4 &	697.5  \\
1099 &	1889 &	965 &	10100 &	600 &	496 &	10720 &	450 &	10700 &	242 &	604.8 &	11745 &	1847.5 \\
1608 &	2869	 & 640 &	10300  & 500 &	1176 &	4380 &	355 &	13100 &	467 &	665.7 &	2826 &	840  \\
\hline
1465.5 &	1181.3 &	734.5 &	14307.8 &	872.5 &	669.3 &	5661 &	471 &	14870 &	384.5 &	600.6 &	4229.6 &	1042.4  \\ 
\hline
397.5 &	856.2 &	129.8 &	7919.3 &	600.1 &	220.6 &	3289.3 &	226.2 &	5964.3 &	170.2  & 89.5 &	2540.6 & 828.4 \\
\hline
\end{tabular}
\vspace*{-0mm}
\caption{CURL implemented on top of Efficient Rainbow - Scores reported for $20$ random seeds for each of the above games, with the last two rows being the mean and standard deviation across the runs.}
\label{tab:atari1}
\vspace{-2mm}
\end{table*}


\begin{table*}[!ht]
\centering
\small
\setlength\tabcolsep{2.5pt}
\begin{tabular}{c|c|c|c|c|c|c|c|c|c|c|c|c}
\hline
UpNDown &	Hero &	CrazyClimber &	ChopperComm. &	DemonAttack &	Amidar &	Alien &	BankHeist &	Breakout &	Freeway &	Pong &	PrivateEye &	Boxing \\
\hline
3529 & 8747.5 &	19090 &	560 &	611.5 &	150.9 &	616 &	95 &	3.6 &	29.2 &	-19.3 &	100 &	-0.5 \\
772 & 3026 &	8290 &	1530 &	707.5 &	131.2 &	923 &	184	& 5 &	25.4 &	-16.9 &	100 &	-11.4 \\
5972 & 7146 &	12160 &	1390 &	843.5 &	141.5 &	467 &	75 &	3.2 &	27.6 &	-12 &	100 &	4 \\
2793 &	7686 &	8920 &	1100 &	330.5 &	133.7 &	441 &	232 &	5.1 &	28.6 &	-19.6 &	100 &	3.6 \\
3546 &	7335 &	11360 &	500 &	759 &	157.1 &	716 &	187 &	2.9 &	22.8 &	-17.8 &	1357.4 &	6.2 \\
4552 &	7325 &	4110 &	990 &	940 &	125.4 &	453 &	367 &	6.3 &	29.6 &	-18.9 &	100 &	5 \\
2972 &	7275.5 &	9460 &	780 &	1136 &	183.2 &	273 &	186 &	5.9 &	23.3 &	-15.9 &	0 &	-1.7 \\
2865 &	3115 &	20630 &	1180 &	758 &	153.6 &	540 &	68 &	2.6 &	27.6 &	-15.2 &	100 &	0.1 \\
3098 &	7424 &	6780 &	1380 &	772.5 &	127.8 &	499 &	60 &	5.9 &	26.1 &	-18.7 &	100 &	3.5 \\
1953 &	7475 &	13570 &	970 &	820 &	149.4 &	475 &	123 &	4.3 &	28.3 &	-13.3 &	100 &	-0.5 \\
1467 &	3135 &	11890 &	1200 &	784 &	125.7 &	553 &	72 &	3.2 &	21.8 &	-17.2 &	1510 &	-22.1 \\
2912 &	5060.5 &	9160 &	1130 &	1080 &	130.4 &	446 &	53 &	4.8 &	21.8 &	-20.1 &	100 &	-1.8 \\
4123 &	4409 &	10960 &	1380 &	847 &	133 &	533 &	68 &	6.3 &	28.9 &	-16.5 &	100 &	1.6 \\
2334 &	6979 &	17360 &	1230 &	771.5 &	140.5 &	968 &	36 &	7.3 &	28.2 &	-14.9 &	100 &	3.6 \\
2605 &	4159 &	8930 &	1350 &	907.5 &	133.8 &	499 &	53 &	4.8 &	28.3 &	-19.3 &	100 &	-17.6 \\
2432 &	7560 &	11510 &	1080 &	1095.5 &	191.8 &	523 &	105 &	3.7 &	26.8 &	-15.6 &	0	 & 21.7 \\
3826 &	8587 &	22690 &	1210 &	700 &	115.5 &	616	 &276	 &6.6 &	27.5	 &-21	 &100	 &2 \\
3052 &	4683.5 &	8120 &	840 &	803.5 &	164 &	475 &	69	 &5.5 &	26.5 &	-10.5 &	0	 &5.9 \\
3131 &	7317 &	13500 &	730 &	818 &	131.7 &	525 &	50	 & 4.3	 & 26.8 &	-13.3 &	100 &	18.7 \\
1169 &	7141 &	14440 &	640 &	866 &	122.4 &	622 &	273 &	6.2 &	28.6 &	-13.1 &	100 &	3.7 \\
\hline
2955.2 & 6279.3 &	12146.5 &	1058.5 &	817.6 &	142.1 &	558.2 &	131.6 &	4.9 &	26.7 &	-16.5 &	218.4 & 1.2 \\
\hline
1181.1 & 1871.5 &	4765.6 &	299.1 &	176.6 &	20.0 &	160.3 &	94.4 &	1.4 &	2.4 & 2.9 & 417.9 &  10.0 \\
\hline
\end{tabular}
\vspace*{-0mm}
\caption{CURL implemented on top of Efficient Rainbow - Scores reported for $20$ random seeds for each of the above games, with the last two rows being the mean and standard deviation across the runs.}
\label{tab:atari2}
\vspace{-2mm}
\end{table*}